\documentclass[11pt]{article}
\usepackage[margin=1in]{geometry}
\usepackage{titlesec}
\usepackage{url}
\usepackage{hyperref}
\usepackage{amsmath}
\usepackage{enumitem} % in preamble
\setlist[itemize]{leftmargin=*, nosep, topsep=0pt, partopsep=0pt}

\newcommand{\memoheader}[4]{
  \noindent
  \textbf{From:} #2\\ 
  \textbf{To:} #1\\
  \textbf{Date:} #3\\
  \textbf{Subject:} #4\\[1ex]\hrule\vspace{1.5ex}
}

\begin{document}
\memoheader{Prof. Barbara Linke}
            {Boston Statics}
            {\today}
            {EME 150B: Project HW 1}

\noindent Given:
\begin{itemize}
    \item $H_{motor} = 1kW$
    \item $n_{motor} = 1000 rpm$
    \item Gear selection:
    \begin{itemize}
        \item 2 gears with 10 teeth, diameter 50 mm
        \item 2 gears with 20 teeth, diameter 100 mm
        \item 2 gears with 30 teeth, diameter 150 mm
        \item 2 gears with 40 teeth, diameter 200 mm
        \item 2 gears with 50 teeth, diameter 250 mm
    \end{itemize}
\end{itemize}

\section*{Subtask 1}
\begin{itemize}
    \item Let radius of teacup platform be: $R_t$
    \item Let mass of child be: $m_c$
    \item let angular speed of platform be: $\omega_t=n_t=1/6$ [rev/sec] 
\end{itemize}

Calculate the forces and corresponding torques, at least for the four load cases:
\begin{enumerate}
    \item momentum of one child in constant motion:
    \begin{align*}
        m_cV = m_c(2\pi n_t R_t)
    \end{align*}
    Torque of one child in constant motion given angular acceleration of $\dot{n_t}$:
    \begin{align*}
        T_t = m_c \dot{n_t} R_t 
    \end{align*}
    \item acceleration/deceleration within one revolution = 6 seconds of one child:
    \begin{align*}
      \dot{n_t}=\frac{T_t}{m_c R_t}
    \end{align*}
    \item momentum of 4 children in constant motion
    \begin{align*}
        m_cV = 4m_c(2\pi n_t R_t)
    \end{align*}
\end{enumerate}

$$ratio=\frac{(10) (10) (20) (30)}{(50) (50) (40) (40) (30)}$$







\end{document}